%% ===== SECTION 5: RESULTS =====
\section{Results}
\label{sec:results}

We organize findings around three diagnostic questions: Is the gap domain-specific? (\S\ref{sec:domain_specific}). Does adaptation compose with prompting? (\S\ref{sec:composability}). How does gap severity vary across benchmarks? (\S\ref{sec:structural_variation}).
Because the preserved artifact coverage differs across benchmarks, tables in this section report the mean over the preserved runs available for each exact model--benchmark--method cell rather than assuming uniform seed coverage everywhere.
Table~\ref{tab:summary} provides a cross-model overview.

\begin{table}[t]
    \caption{\textbf{Cross-model summary.} Accuracy~(\%) from the preserved runs used in this draft.
    ``Best prompt'' = strongest frozen prompting method.
    ``Best overall'' = strongest configuration including BSTTT.
    $\Delta$: best overall minus frozen.}
    \label{tab:summary}
    \centering
    \small
    \begin{tabular}{@{}llcccc@{}}
        \toprule
        \textbf{Model} & \textbf{Condition} & \textbf{SimpleToM} & \textbf{Hi-ToM} & \textbf{OpenToM} & \textbf{DynToM} \\
        \midrule
        \multirow{4}{*}{Llama-8B}
            & Frozen         & 55.6 & 41.4 & 42.5 & 19.0 \\
            & Best prompt    & 70.0 & 47.5 & 48.8 & 19.5 \\
            & BSTTT-AR       & 61.1 & 54.4 & 43.1 & 18.2 \\
            & Best overall   & \textbf{83.3} & \textbf{55.3} & \textbf{50.0} & 22.8 \\
        \cmidrule(l){2-6}
            & $\Delta$       & +27.7 & +13.9 & +7.5 & +3.8 \\
        \midrule
        \multirow{4}{*}{Qwen-7B}
            & Frozen         & 41.7 & 40.7 & 40.5 & 32.5 \\
            & Best prompt    & 62.7 & 43.8 & 45.6 & 33.0 \\
            & BSTTT-AR       & 61.6 & 40.7 & 40.5 & --- \\
            & Best overall   & \textbf{73.3} & \textbf{53.7} & \textbf{46.7} & 35.0 \\
        \cmidrule(l){2-6}
            & $\Delta$       & +31.6 & +13.0 & +6.2 & +2.5 \\
        \midrule
        \multirow{4}{*}{GPT-oss-20B}
            & Frozen         & 66.7 & 32.6 & 35.6 & 15.5 \\
            & Best prompt    & 66.7 & 34.8 & 35.8 & 16.0 \\
            & BSTTT-AR       & 66.7 & 47.3 & 36.0 & --- \\
            & Best overall   & 66.7 & \textbf{47.3} & \textbf{36.3} & 16.0 \\
        \cmidrule(l){2-6}
            & $\Delta$       & +0.0 & +14.7 & +0.7 & +0.5 \\
        \bottomrule
    \end{tabular}
\end{table}


\subsection{Finding 1: The Gap Reflects Latent Competence, Not Absent Knowledge}
\label{sec:domain_specific}

The central diagnostic question is whether the gap reflects a capacity limitation (\textbf{H1}) or recoverable latent competence (\textbf{H2}).
Table~\ref{tab:ar_vs_ntl} presents the critical contrast.

\begin{table}[t]
    \caption{\textbf{The diagnostic contrast: AR vs.\ NTL.}
    Both conditions use identical LoRA configuration, support data, and gradient budget.
    Cells report the mean over the preserved runs available for that exact setting.
    $\Delta$(AR$-$NTL) isolates the contribution of ToM-aligned supervision.}
    \label{tab:ar_vs_ntl}
    \centering
    \small
    \begin{tabular}{@{}llcccr@{}}
        \toprule
        \textbf{Model} & \textbf{Benchmark} & \textbf{Frozen} & \textbf{NTL} & \textbf{AR} & $\boldsymbol{\Delta}$\textbf{(AR$-$NTL)} \\
        \midrule
        \multirow{3}{*}{Llama-8B}
            & SimpleToM  & 55.6 & 55.6 & 61.1 & +5.5 \\
            & Hi-ToM     & 41.4 & 41.4 & 54.4 & \textbf{+13.0} \\
            & OpenToM    & 42.5 & 42.7 & 43.1 & +0.4 \\
        \midrule
        \multirow{3}{*}{Qwen-7B}
            & SimpleToM  & 41.7 & 45.3 & 61.6 & \textbf{+16.3} \\
            & Hi-ToM     & 40.7 & 40.7 & 40.7 & +0.0 \\
            & OpenToM    & 40.5 & 40.5 & 40.5 & +0.0 \\
        \midrule
        \multirow{3}{*}{GPT-oss-20B}
            & SimpleToM  & 66.7 & 66.7 & 66.7 & +0.0 \\
            & Hi-ToM     & 32.6 & 31.6 & 47.3 & \textbf{+15.7} \\
            & OpenToM    & 35.6 & 35.1 & 36.0 & +0.9 \\
        \bottomrule
    \end{tabular}
\end{table}

The results support \textbf{H2} on the benchmarks where AR substantially outperforms NTL.
On Hi-ToM, BSTTT-AR recovers +13.0~pp (Llama) and +15.7~pp (GPT-oss) over frozen inference, while BSTTT-NTL provides \emph{zero improvement}---despite using the same model, adapters, learning rate, and gradient steps.
On SimpleToM, Qwen shows the same pattern: AR yields +16.3~pp while NTL yields +3.6~pp.
The AR--NTL contrast rules out generic explanations: the only difference is \emph{what is predicted}. The gap is domain-specific.

\paragraph{Sub-question analysis.}
Table~\ref{tab:simpletom_subtypes} shows BSTTT-AR selectively improves \emph{applied} accuracy while leaving \emph{explicit} accuracy unchanged.

\begin{table}[t]
    \caption{\textbf{SimpleToM per-question-type accuracy, before and after BSTTT-AR.}
    BSTTT selectively improves behavior prediction (applied) while mental-state accuracy (explicit) is stable.
    Values use the preserved multi-seed SimpleToM runs available for each model.}
    \label{tab:simpletom_subtypes}
    \centering
    \small
    \begin{tabular}{@{}l ccc ccc ccc @{}}
        \toprule
        & \multicolumn{3}{c}{\textbf{Qwen-7B}} & \multicolumn{3}{c}{\textbf{Llama-8B}} & \multicolumn{3}{c}{\textbf{GPT-oss-20B}} \\
        \cmidrule(lr){2-4} \cmidrule(lr){5-7} \cmidrule(lr){8-10}
        \textbf{Question type} & Bef. & Aft. & $\Delta$ & Bef. & Aft. & $\Delta$ & Bef. & Aft. & $\Delta$ \\
        \midrule
        Mental-state (\emph{expl.}) & 54.4 & 54.4 & 0.0 & 100 & 100 & 0.0 & 100 & 100 & 0.0 \\
        Behavior (\emph{appl.})     & 0.0 & \textbf{54.4} & \textbf{+54.4} & 26.7 & 43.3 & +16.7 & 100 & 100 & 0.0 \\
        Judgment (\emph{appl.})     & 76.0 & 76.0 & 0.0 & 40.0 & 40.0 & 0.0 & 0.0 & 0.0 & 0.0 \\
        \midrule
        Average & 43.5 & 61.6 & +18.1 & 55.6 & 61.1 & +5.6 & 66.7 & 66.7 & 0.0 \\
        \bottomrule
    \end{tabular}
\end{table}

The Qwen result is particularly striking: the model achieves \emph{zero} correct behavior predictions under frozen inference despite correctly identifying the mental state 54\% of the time.
After three gradient steps, behavior accuracy matches mental-state accuracy at 54.4\%---completely closing the gap.
Llama shows a narrower but still targeted effect: BSTTT improves behavior prediction from 26.7\% to 43.3\% while leaving both mental-state and judgment accuracy unchanged.
GPT-oss shows a ceiling effect on SimpleToM (already 100\% on behavior), with its gap manifesting exclusively in judgment questions (0\%), which BSTTT does not improve.

\begin{tcolorbox}[colback=gray!5, colframe=gray!50, boxrule=0.3pt, left=4pt, right=4pt, top=3pt, bottom=3pt]
\textbf{Evaluation implication:}
Aggregate accuracy hides systematic dissociations.
Qwen's 41.7\% aggregate obscures 54\% explicit vs.\ 0\% applied accuracy---the wrong conclusion about partial understanding, when the truth is complete recognition with complete application failure.
\end{tcolorbox}


\subsection{Finding 2: BSTTT Composes with Prompting}
\label{sec:composability}

BSTTT and prompting are fundamentally different interventions: one adapts model weights, the other modifies model inputs.
Table~\ref{tab:composability} tests whether their gains are additive.

\begin{table}[t]
    \caption{\textbf{Composability: BSTTT composes with prompting.}
    Accuracy~(\%) from the preserved runs used in this draft.
    ``BSTTT'' = BSTTT-AR with vanilla prompting.
    ``BSTTT+Prompt'' = BSTTT-AR composed with the best prompting strategy.
    $\Delta$: BSTTT+Prompt minus Frozen.}
    \label{tab:composability}
    \centering
    \small
    \begin{tabular}{@{}llcccr@{}}
        \toprule
        \textbf{Model} & \textbf{Benchmark} & \textbf{Frozen} & \textbf{BSTTT} & \textbf{BSTTT+Prompt} & $\boldsymbol{\Delta}$ \\
        \midrule
        \multirow{4}{*}{Llama-8B}
            & SimpleToM  & 55.6 & 61.1  & 83.3  & \textbf{+27.7} \\
            & Hi-ToM     & 41.4 & 54.4  & 55.3  & +13.9 \\
            & OpenToM    & 42.5 & 43.1  & 50.0  & +7.5 \\
            & DynToM     & 19.0 & 18.2  & 22.8  & +3.8 \\
        \midrule
        \multirow{4}{*}{Qwen-7B}
            & SimpleToM  & 41.7 & 61.6  & 73.3  & \textbf{+31.6} \\
            & Hi-ToM     & 40.7 & 40.7  & 53.7  & +13.0 \\
            & OpenToM    & 40.5 & 40.5  & 46.7  & +6.2 \\
            & DynToM     & 32.5 & 32.5  & 35.0  & +2.5 \\
        \midrule
        \multirow{4}{*}{GPT-oss-20B}
            & SimpleToM  & 66.7 & 66.7  & 66.7  & +0.0 \\
            & Hi-ToM     & 32.6 & 47.3  & 46.5  & +13.9 \\
            & OpenToM    & 35.6 & 36.0  & 36.3  & +0.7 \\
            & DynToM     & 15.5 & 16.5  & 16.0  & +0.5 \\
        \bottomrule
    \end{tabular}
\end{table}

Three patterns emerge.

\paragraph{Prompting amplifies adaptation.}
On 9 of 12 model--benchmark cells, the combined intervention exceeds BSTTT alone---often substantially.
The strongest amplification occurs on SimpleToM, where Llama improves from 61.1\% (BSTTT alone) to 83.3\% (BSTTT+Prompt), and Qwen from 61.6\% to 73.3\%.
On OpenToM, Llama's BSTTT-alone gain is marginal (+0.6~pp), but BSTTT+Prompt achieves +7.5~pp, indicating that prompting structures the input in a way that makes the adapted weights more effective.

\paragraph{Prompting can unlock otherwise-inert adaptation.}
The Qwen results on Hi-ToM and OpenToM are particularly revealing.
BSTTT-AR \emph{alone} produces zero improvement on either benchmark (identical to Frozen).
Yet when paired with structured prompting, accuracy jumps to 53.7\% on Hi-ToM (+13.0~pp) and 46.7\% on OpenToM (+6.2~pp).
This is not simply additive---BSTTT alone contributes nothing, yet the combination far exceeds prompting alone (Scratchpad achieves 43.8\% on Hi-ToM).
The prompt provides a representational scaffold---surfacing belief-relevant features---that gives the LoRA adaptation meaningful gradient signal.

\paragraph{Diminishing returns at the ceiling.}
When one intervention already captures most of the recoverable gap, the other adds little.
On GPT-oss Hi-ToM, BSTTT-AR alone achieves +14.7~pp; adding prompting provides no further benefit and marginally decreases accuracy (47.3\%~$\rightarrow$~46.5\%), suggesting the scaffolding can interfere with adapted representations.

\begin{tcolorbox}[colback=gray!5, colframe=gray!50, boxrule=0.3pt, left=4pt, right=4pt, top=3pt, bottom=3pt]
\textbf{Evaluation implication:}
Testing interventions in isolation underestimates joint potential.
On Qwen Hi-ToM, BSTTT alone shows zero benefit; Scratchpad alone gives +3.1~pp.
Only the combined evaluation reveals 13~pp of accessible latent competence.
\end{tcolorbox}


\subsection{Finding 3: Temporal Belief Tracking Is a Genuine Capacity Bottleneck}
\label{sec:structural_variation}

The magnitude of BSTTT's diagnostic signal varies substantially across benchmarks.

\paragraph{Hi-ToM and SimpleToM: largest gains.}
Hi-ToM shows the most consistent gains across all three models (+13.0 to +14.7~pp for BSTTT-AR alone), while SimpleToM provides the cleanest diagnostic signal (0\%$\rightarrow$54\% on Qwen).

\paragraph{OpenToM: gains contingent on prompting.}
BSTTT-AR alone provides only marginal improvement on OpenToM.
However, combining BSTTT with structured prompting yields meaningful gains (up to +7.5~pp).
OpenToM's multihop questions require chaining mental-state attributions---BSTTT alone cannot compensate for missing intermediate reasoning steps.

\paragraph{DynToM: temporal reasoning resists everything.}
DynToM yields the smallest improvements across \emph{all} methods.
The best configuration achieves only +3.8~pp (Llama), +2.5~pp (Qwen), and +0.5~pp (GPT-oss).
Even Scratchpad-Oracle---providing ground-truth mental-state tables---reaches only 35.0\% on Qwen, far below ceiling.
This indicates that DynToM's difficulty is not reducible to the explicit-to-applied gap: the bottleneck is \emph{temporal reasoning over evolving belief states}, a capacity that neither adaptation nor prompting can reliably provide.

\begin{tcolorbox}[colback=gray!5, colframe=gray!50, boxrule=0.3pt, left=4pt, right=4pt, top=3pt, bottom=3pt]
\textbf{Evaluation implication:}
Temporal and static ToM probe fundamentally different competencies.
Benchmarks testing only static scenarios overestimate social reasoning capability.
DynToM-style temporal tasks should be evaluated as a separate axis.
\end{tcolorbox}


\subsection{Error Analysis}
\label{sec:errors}

We examine how BSTTT-AR redistributes errors across benchmarks by comparing per-instance predictions before and after adaptation (GPT-oss-20B, seed~42).

\paragraph{Hi-ToM: broad, net-positive correction.}
Of 1{,}200 items, BSTTT-AR flips 263 incorrect frozen predictions to correct (21.9\%) while introducing 96 new errors (8.0\%)---a 2.7:1 correction-to-regression ratio.
The remaining 546 items (45.5\%) are incorrect under both conditions.
This indicates that gains are broadly distributed across scenarios, not concentrated on a small subset of easy items.

\paragraph{OpenToM: near-zero churn.}
Of 2{,}000 items, BSTTT-AR corrects only 19 items (0.9\%) and breaks 9 (0.5\%), leaving 1{,}268 items (63.4\%) persistently wrong.
This low churn confirms quantitatively what the aggregate numbers suggest: BSTTT-AR alone barely touches OpenToM.
The decomposition by question type is revealing: attitude questions (explicit) show exactly zero change (35.1\% frozen $=$ 35.1\% after BSTTT), while multihop questions (applied) show a marginal +0.7~pp gain.
The persistent errors on multihop questions require multi-step reasoning chains that episodic LoRA adaptation cannot construct \emph{de novo}---it can only strengthen existing pathways.

\paragraph{SimpleToM: asymmetric across question types.}
BSTTT-AR's corrections are concentrated entirely on behavior questions (the applied component), while mental-state questions (explicit) and judgment questions exhibit zero change (Table~\ref{tab:simpletom_subtypes}).
This asymmetry is diagnostic: it shows that BSTTT activates a specific latent capability (belief-to-behavior mapping) rather than providing a generic accuracy boost.
The persistent failure on judgment questions---even after successful behavior recovery---suggests that \emph{evaluative} social reasoning (``Was this action reasonable?'') involves qualitatively different mechanisms than \emph{predictive} reasoning (``What will they do?'').

\paragraph{Taxonomy of persistent failures.}
Qualitative review of items incorrect across all methods and seeds reveals three dominant patterns:
(i)~\emph{support--query mismatch}: support examples involve agents or scenarios too dissimilar from the query for adaptation to transfer (most common on OpenToM, where narrative complexity varies widely);
(ii)~\emph{deep compositional failure}: the model lacks capacity for the required reasoning depth regardless of adaptation (e.g., 4th-order recursive belief in Hi-ToM);
(iii)~\emph{world-knowledge conflicts}: the model defaults to its prior world knowledge rather than reasoning from the agent's perspective---for example, predicting that a character will \emph{not} buy expired yogurt even when the character is unaware it is expired (common on SimpleToM judgment questions).
These represent genuine capacity limitations (\textbf{H1}) coexisting with the recoverable component (\textbf{H2}), demonstrating that the explicit-to-applied gap has both addressable and fundamental components.


%% ===== SECTION 6: RECOMMENDATIONS =====
\section{Recommendations for Theory of Mind Evaluation}
\label{sec:recommendations}

Our findings motivate four concrete recommendations.

\paragraph{R1: Decompose aggregate scores.}
Separately report accuracy on explicit and applied questions wherever benchmark structure permits.
Aggregate accuracy can inflate apparent competence by 5--54~pp by averaging strong recognition with weak application.

\paragraph{R2: Distinguish temporal from static reasoning.}
Temporal belief tracking is qualitatively harder than static inference and resists all interventions tested.
ToM evaluations should treat temporal reasoning as a separate axis, not a harder version of the same skill.

\paragraph{R3: Include adaptation-based diagnostics.}
Standard evaluation provides a lower bound on capability but cannot determine \emph{why} a model fails.
Controlled adaptation experiments estimate the upper bound latent in model parameters, informing whether future work should target better inference or better training.

\paragraph{R4: Test composed interventions.}
Evaluating prompting and adaptation in isolation underestimates joint potential.
Combined interventions can exceed either component by 10+~pp.


%% ===== SECTION 7: RELATED WORK =====
\section{Related Work}
\label{sec:related}

\paragraph{ToM evaluation in LLMs.}
Early false-belief task adaptations reported near-human LLM performance~\citep{kosinski2024evaluating}, but subsequent work revealed failures under trivial perturbations~\citep{ullman2023large} and adversarial probes~\citep{shapira2024clever}, with~\citet{sap2022neural} arguing apparent ToM reflects surface heuristics.
The benchmark landscape now includes SimpleToM~\citep{gu2024simpletom}, Hi-ToM~\citep{wu-etal-2023-hi}, OpenToM~\citep{xu-etal-2024-opentom}, DynToM~\citep{xiao-etal-2025-towards}, and FANToM~\citep{kim-etal-2023-fantom}.
Our work differs in diagnostic aim: we investigate \emph{why} models fail to apply recognized mental states and whether the failure is recoverable.

\paragraph{Prompting for social reasoning.}
CoT improves explicit ToM but inconsistently affects applied prediction~\citep{moghaddam2023boosting}.
SimToM~\citep{wilf2024think} decomposes perspective-taking; SymbolicToM~\citep{sclar-etal-2023-minding} scaffolds via belief-state graphs.
These modify inputs; BSTTT modifies weights. We show the two compose---a finding prompting-only evaluations cannot reveal.

\paragraph{Test-time training.}
TTT adapts parameters at evaluation time~\citep{Sun2019TestTimeTW, NEURIPS2021_b618c321}.
Recent work applies TTT to LMs~\citep{feng2026inplace}.
Our contribution is the contrastive AR--NTL design that transforms TTT into a \emph{diagnostic instrument}---a use not explored in prior work.


%% ===== SECTION 8: LIMITATIONS =====
\section{Limitations and Broader Impact}
\label{sec:limitations}

\paragraph{Scale.} We evaluate models at 7--20B.
The gap may narrow at larger scales; our protocol is scale-agnostic.

\paragraph{Format.} All benchmarks use multiple-choice QA.
Open-ended generation may reveal different patterns.

\paragraph{Indirect measurement.} Hi-ToM contains only applied questions in our operationalization.
Our Hi-ToM conclusions therefore concern recoverability under adaptation rather than a directly observed explicit-to-applied gap within that benchmark.

\paragraph{Overhead.} BSTTT approximately triples inference cost ($K{=}3$ gradient steps).
Acceptable for diagnostics; deployment requires amortization.

\paragraph{Temporal mismatch.} BSTTT's episodic reset is poorly matched to temporal tasks like DynToM.
Continual adaptation may be more appropriate and warrants investigation.

\paragraph{Artifact coverage.} The preserved artifact bundle used in this draft is stronger for some benchmark/model combinations than others, especially on SimpleToM and DynToM.
We therefore report preserved multi-seed means rather than a perfectly uniform seed protocol for every cell, and future work should rerun the full matrix under a single standardized protocol.

\paragraph{Broader impact.} Our finding that applied ToM is often \emph{latent} rather than absent has implications for deployment in social contexts.
A model that \emph{can} reason about mental states but does not \emph{reliably} do so may present deceptive capability profiles.
We encourage decomposed scoring to avoid overestimating deployed model competence.


%% ===== SECTION 9: CONCLUSION =====
\section{Conclusion}
\label{sec:conclusion}

We introduced the explicit-to-applied Theory of Mind gap and a diagnostic evaluation protocol (BSTTT) to characterize its nature.
Across four benchmarks and three models, we showed that the gap is largely recoverable through domain-aligned adaptation (supporting the latent-competence hypothesis), that weight-level and input-level interventions compose, and that temporal belief tracking remains a genuine bottleneck.
Standard aggregate ToM scores can overstate applied competence while obscuring recoverable latent competence.
We recommend that future evaluations decompose scores, include temporal reasoning as a separate axis, and use adaptation-based diagnostics to bound the gap between what models \emph{do} and what their parameters appear able to support under controlled intervention.
An anonymized code and artifact bundle accompanies the submission.


\bibliography{references}
\bibliographystyle{plainnat}

\newpage
\appendix
\section{Full Per-Model Results}
\label{app:full_results}

Table~\ref{tab:modelwise_llama} through Table~\ref{tab:modelwise_gptoss} provide the consolidated per-model results recovered from the preserved artifact bundle and verification notes used in this draft.
% Auto-generated consolidated model-wise tables with DynToM included.
% Source bundle: recovered commits + synced local summaries.

\begin{table}[t]
\caption{\textbf{Llama-3.1-8B results} (accuracy, \%).}
\label{tab:modelwise_llama}
\centering
\small
\begin{tabular}{lcccc}
\toprule
\textbf{Method} & \textbf{SimpleToM} & \textbf{DynToM} & \textbf{Hi-ToM} & \textbf{OpenToM} \\
\midrule
Frozen & 55.56 & 19.00 & 41.42 & 42.45 \\
CoT & 62.22 & --- & 38.75 & 41.55 \\
Vanilla & --- & --- & --- & --- \\
Scratchpad-Frozen & 70.00 & 19.50 & 43.00 & 45.80 \\
SimToM & --- & --- & 41.83 & 48.10 \\
SymbolicToM & --- & --- & 47.50 & 48.75 \\
BSTTT-NTL & 55.56 & --- & 41.39 & 42.67 \\
BSTTT-AR & 61.11 & --- & 54.36 & 43.13 \\
Scratchpad+BSTTT-AR & --- & --- & 53.72 & --- \\
SimToM+BSTTT-AR & --- & --- & 53.86 & 47.65 \\
SymbolicToM+BSTTT-AR & --- & --- & 55.25 & 50.02 \\
Scratchpad-TTT (DynToM) & --- & 22.83 & --- & --- \\
Scratchpad-Oracle (DynToM) & --- & 20.50 & --- & --- \\
Hierarchical-TTT (DynToM) & --- & 16.99 & --- & --- \\
\bottomrule
\end{tabular}
\vspace{2pt}
\caption*{\footnotesize DynToM summaries in recovered artifacts do not always include an explicit model tag; these DynToM values are from the same recovered run batch as the Llama results (commit \texttt{a5d8721d}).}
\end{table}

\begin{table}[t]
\caption{\textbf{Qwen-2.5-7B results} (accuracy, \%).}
\label{tab:modelwise_qwen}
\centering
\small
\begin{tabular}{lcccc}
\toprule
\textbf{Method} & \textbf{SimpleToM} & \textbf{DynToM} & \textbf{Hi-ToM} & \textbf{OpenToM} \\
\midrule
Frozen & 41.67 & 32.50 & 40.67 & 40.50 \\
CoT & 52.67 & --- & 42.58 & 41.00 \\
Vanilla & 45.33 & --- & --- & --- \\
Scratchpad-Frozen & 62.67 & 33.00 & 43.83 & 40.40 \\
SimToM & --- & --- & 38.83 & 45.60 \\
SymbolicToM & --- & --- & 43.00 & 42.70 \\
BSTTT-NTL & 45.33 & 32.50 & 40.75 & 40.55 \\
BSTTT-AR & 61.60 & --- & 40.75 & 40.55 \\
Scratchpad+BSTTT-AR & --- & --- & 53.75 & 36.90 \\
SimToM+BSTTT-AR & --- & --- & 49.06 & 46.70 \\
SymbolicToM+BSTTT-AR & --- & --- & --- & 39.28 \\
Scratchpad-TTT (DynToM) & --- & 33.50 & --- & --- \\
Scratchpad-Oracle (DynToM) & --- & 35.00 & --- & --- \\
Hierarchical-TTT (DynToM) & --- & 33.33 & --- & --- \\
\bottomrule
\end{tabular}
\vspace{2pt}
\caption*{\footnotesize DynToM Qwen values are from the saved 200-question slice diagnostics (same settings as Direction-2 runs), as documented in recovered logs/tables.}
\end{table}

\begin{table}[t]
\caption{\textbf{GPT-oss-20B results} (accuracy, \%).}
\label{tab:modelwise_gptoss}
\centering
\small
\begin{tabular}{lcccc}
\toprule
\textbf{Method} & \textbf{SimpleToM} & \textbf{DynToM} & \textbf{Hi-ToM} & \textbf{OpenToM} \\
\midrule
Frozen & 66.67 & 16.00 & 32.58 & 35.65 \\
CoT & --- & --- & 32.67 & 31.50 \\
Vanilla & --- & --- & --- & --- \\
Scratchpad-Frozen & --- & --- & 34.83 & 35.80 \\
SimToM & --- & --- & 31.72 & 34.40 \\
SymbolicToM & --- & --- & 25.67 & 28.85 \\
BSTTT-NTL & --- & --- & 31.58 & 35.13 \\
BSTTT-AR & --- & --- & 47.33 & 35.97 \\
Scratchpad+BSTTT-AR & --- & --- & 46.50 & 36.25 \\
SimToM+BSTTT-AR & --- & --- & 41.19 & 34.90 \\
SymbolicToM+BSTTT-AR & --- & --- & 32.50 & 27.47 \\
Scratchpad-TTT (DynToM) & --- & --- & --- & --- \\
Scratchpad-Oracle (DynToM) & --- & --- & --- & --- \\
Hierarchical-TTT (DynToM) & --- & --- & --- & --- \\
\bottomrule
\end{tabular}
\vspace{2pt}
\caption*{\footnotesize GPT-oss-20B DynToM currently has only a recovered Frozen run (seed 42) in commit \texttt{0e1cc228}.}
\end{table}


\section{Hyperparameter Sensitivity}
\label{app:sensitivity}

Our primary goal is diagnostic control rather than benchmark-specific optimization, so we keep the adaptation configuration fixed across AR and NTL comparisons: LoRA rank $r{=}8$, $\alpha{=}16$, top 25\% of layers, AdamW with learning rate $\eta{=}10^{-4}$, and $K{=}3$ adaptation steps.
We did not run a fully standardized hyperparameter sweep across every benchmark/model pair in the preserved artifact set; instead, we use this single configuration to make the contrast between aligned and generic supervision interpretable.
This choice should be read as a limitation of the present study rather than evidence that the reported settings are globally optimal.
